\begin{table}[htbp]\centering
\small
\centering
	\caption{\textsc{List of STEM majors}}
 \label{tab: listSTEMmajors}
 \begin{threeparttable}
    \begin{tabular*}{\textwidth}{@{\extracolsep{\fill}}l@{}} 
 \toprule
        STEM majors\\ 
\midrule
        
        Biological Sciences \\ 
        Physical Sciences \\
        Natural Sciences, Mathematics and Statistics \\ 
        Industry and Production \\ 
        Engineering and Related Professions \\         Environmental Studies \\ 
                Forestry, Agriculture, Fisheries \\
        Health \\ 
        Information and Communication Technology \\ 
        Veterinary Medicine \\ 
\bottomrule
    \end{tabular*}
\begin{tablenotes}\singlespacing
 \item   \scriptsize \noindent \textsc{ Note.--}  Source: \cite{tincani_miglino_university_major_classifications_2026}. The major categorization corresponds to the subarea categorization established by UNESCO in the CINE-F classification defined in 2013 and being used by the OECD since 2016 \citep{listmajors}. Data provided by the Chilean Ministry of Education \citep{mineduc2023sistema}. The distinction between STEM and Non-STEM majors follows the definition of STEM disciplines provided by the \citet{STEMclassification}. 
 \end{tablenotes}
  \end{threeparttable}
\end{table}